% HD-S43. GPT-6 Astra Pro /root/sixfolds; inherited effort not exposed.
% 21 November 2026. Actual tensor quotient and universal period bound.
\subsection{HD-S43: exclusion of the multiplicity-two tensor case}
\label{sec:two-multiplicity-tensor-exclusion}

For every imaginary quadratic exact coefficient field \(K\),
the multiplicity-two case of HD-S42 has period rank at most
eight for its natural signature-\((3,3)\) sixfold factor.
The statement concerns the prescribed actual primitive corner
and restricted theta polarization. Its proof uses an actual
rank-two coefficient system on a quotient curve and an
explicit complementary Jacobian projector. No classification
of finite subgroups of \(GL_2(K)\), determinant-lattice
assumption, or numerical spectral search is needed.

\paragraph{1. The exact reflection tensor and its high-rank alternatives.}
Suppose, to obtain a contradiction, that the period rank is
nine. The maps (S42.8)--(S42.13) reduce the source to a
faithful coefficient group \(H\) on
\[
             U=W\otimes_K M,\qquad \dim_KW=r,
                   \dim_K M=2,\qquad r\ge12.              \tag{S43.1}
\]
The normal supported subgroup is an actual absolutely
irreducible \(K\)-reflection group \(L\) on \(W\),
acting identically on \(M\). The group \(H\) is absolutely
irreducible, normalizes \(\operatorname{End}_K(W)\otimes I\),
and has exact character field \(K\). Its base is
\(\mathbb P^1\); its inertia tuple has \(b\ge12\)
nonidentity entries and total codimension \(4r+6\).
All these assertions, including \(r\ge12\), are proved
by the actual tensor and quotient maps of HD-S42.

Cohen's reflection classification in this range gives either
the standard \(A_r\) reflection representation, or the
imprimitive family \(G(\ell,e,r)\). The exact primary
locators are Cohen (1976), Theorem(2.4), pp.386--387;
(4.4)(i), p.399; table(4.16), p.412; and Theorem(5.12),
p.435. The nonreal primitive exceptions have rank at most
six and the exceptional real types have rank at most eight.
The signed permutation families are included in the
imprimitive alternative. This is the same high-rank list
and actual standard models used in HD-G30.

In the imprimitive case the trace recovery (G30.8)--(G30.10)
shows that its root \(\zeta_\ell\) lies in the actual
character field of \(L\), hence in \(K\).
Its diagonal subgroup \(D\) is the unique largest abelian
normal subgroup, since \(L/D=S_r\) and \(r\ge5\).
The distinct coordinate characters therefore give the literal
\(K\)-rank-one projectors on \(W\)
\[
       f_j=|D|^{-1}\sum_{d\in D}
                          \lambda_j(d)^{-1}\rho_W(d).
                                                               \tag{S43.2}
\]
The actual \(H\)-normalizer preserves \(L\), hence
\(D\), and permutes these projectors. Thus
\(f_jW\otimes M\) is a transitive \(K\)-rank-two
block system for \(H\). Its block stabilizer is absolutely
irreducible, because a proper invariant subspace would induce
one for \(H\). Orthogonality follows from a diagonal
element acting identically on one axis and nontrivially on
another, as in HD-G30. The actual rank-two induction,
quotient and corner maps of HD-S30 give period rank at most
seven. This excludes the imprimitive alternative.

\paragraph{2. The alternating standard module gives an actual product group.}
Put \(s=r+1\). In the remaining case \(L\simeq S_s\),
and its \(K\)-module is the standard augmentation space
\[
                    W=\{(x_1,\ldots,x_s)\in K^s:
                                             \sum x_i=0\}.
                                                               \tag{S43.3}
\]
This is an actual \(K\)-intertwiner: the finite equations
with the rational standard matrices have a one-dimensional
complex solution space, so Gaussian elimination over \(K\)
gives a nonzero, and hence invertible, solution. This is the
same exact descent as (S41.12).

The direct cubic-Hessian proof (S41.6)--(S41.9) identifies
the full complex normalizer of standard \(A_s\), and
therefore the normalizer of \(S_s\), with scalar multiples
of standard permutations. Apply it to the first factor of
the actual tensor normalizer from (S42.11). Absorb its
scalar into the second factor. Every member of \(H\)
then has the unique literal form
\[
                         P_\pi|_W\otimes B,
                                \qquad B\in GL_2(K).       \tag{S43.4}
\]
Uniqueness follows because a standard permutation is scalar
only when it is the identity. In the original \(K\)-basis,
multiplication by the inverse rational permutation leaves
\(I_W\otimes B\), so \(B\) has entries in \(K\).
Consequently (S43.4) is an actual faithful homomorphism
\(H\to S_s\times GL_2(K)\).

The supported reflection subgroup is the full
\(S_s\times1\). If \(\Gamma\) denotes the second
projection, multiplication of any \((\pi,B)\) by
\((\pi^{-1},1)\) shows that the image is exactly
\[
                              H=S_s\times\Gamma.           \tag{S43.5}
\]
The group \(\Gamma\subset GL_2(K)\) is finite and
absolutely irreducible; a common line would make its tensor
with \(W\) invariant under \(H\). Its exact character
field is \(K\), since the trace on \((1,B)\) is
\(r\operatorname{tr}(B)\). The prescribed invariant
positive coefficient form has the literal product expression
\(h_W\otimes h_M\), with the standard rational form
\(h_W\) on (S43.3). To see this, fix a nonzero standard
vector \(w\) and define
\(h_M(a,b)=h_U(w\otimes a,w\otimes b)/h_W(w,w)\);
Schur averaging over \(S_s\) proves this expression for
all pairs of vectors. Thus the coefficient polarizations
are specified as well as the representations.

\paragraph{3. Exactly three nonidentity rank-two monodromies.}
Write the generating sphere inertia as
\((\pi_j,\gamma_j)\in S_s\times\Gamma\).
The projective image of \(\Gamma\) on \(M\) is
noncyclic, since a finite cyclic projective image has a
common eigenline. Therefore at least three \(\gamma_j\)
are projectively nonidentity. For each such entry, its two
distinct eigenvalues imply fixed dimension at most \(r\)
on \(W\otimes M\); its codimension is at least \(r\).
Every nonidentity tensor has codimension at least two.
Four projectively nonidentity entries would give
\[
                    \sum c_j\ge4r+2(b-4)
                                      \ge4r+16>4r+6.
                                                               \tag{S43.6}
\]
There are exactly three.

There can be no additional nonidentity scalar \(\gamma_j\).
For a scalar root of unity \(a\ne1\) of order \(q\ge2\),
the fixed space on \(K^s\otimes M\) has dimension twice
the number of permutation cycles with lengths divisible
by \(q\), at most \(2s/q\le s\). The constant
\(M\)-summand has no fixed vector, so the same bound
holds on \(W\otimes M\), giving codimension at least
\(2r-s=r-1\). Together with the three projective entries
and the other \(b-4\ge8\) nonidentity entries, this gives
\[
                   \sum c_j\ge3r+(r-1)+16
                                      =4r+15>4r+6.         \tag{S43.7}
\]
Thus precisely three \(\gamma_j\) are nonidentity;
they are projectively nonidentity, generate \(\Gamma\),
and have product one. All other inertia have
\(\gamma_j=1\).

Let \(\mathcal M\) denote the actual finite coefficient
system of rank two on \(\mathbb P^1\) with these three
monodromies in \(\Gamma\). Its invariant and
dual invariant spaces vanish. Therefore
\[
 d_0=\dim_K H^1(\mathbb P^1,j_*\mathcal M)
       =-4+\sum_{j=1}^3\operatorname{rank}(\gamma_j-1)
                         \in\{0,1,2\}.                   \tag{S43.8}
\]
Nonnegativity is the dimension of this actual cohomology
group. Its cover and Hodge structure are constant after
fixing the finite monodromy labels and moving the three
branch values to \(0,1,\infty\).

\paragraph{4. The degree-\(s\) quotient and the exact complementary factor.}
Let \(C\to\mathbb P^1\) be the faithful full cover in
(S43.5), and let \(L_0=S_{s-1}\) stabilize the first
permutation letter. Form the actual quotients
\[
 \begin{gathered}
 X=C/(L_0\times1),\qquad Z=C/(S_s\times1),\qquad
 Y=C/(L_0\times\Gamma),\\
 \alpha:C\to X,\qquad f:X\to Z,\qquad\deg f=s.
 \end{gathered}                                               \tag{S43.9}
\]
The subgroup intersection and product are respectively
\(L_0\times1\) and \(H\), so the map on quotient
fibres identifies \(X\) with the normalization of
\(Y\times_{\mathbb P^1}Z\). Both \(X/Y\) and
\(Z/\mathbb P^1\) have full deck group \(\Gamma\),
and \(f\) is \(\Gamma\)-equivariant. In particular
the pulled-back rank-two system on \(Y\) is absolutely
irreducible.

The permutation coefficient system splits by the literal maps
\[
 \begin{gathered}
 K^s\otimes M=(\boldsymbol1\otimes M)\perp(W\otimes M),\\
 j(a)=(a,\ldots,a),\qquad
 p(a_1,\ldots,a_s)=s^{-1}\sum_i a_i,\qquad pj=1.
 \end{gathered}                                               \tag{S43.10}
\]
The constant copy of \(M\) has form \(s h_M\).
On the cover fibres these are precisely the direct-image and trace
maps of \(Y\to\mathbb P^1\); a permutation cycle of
length \(\ell\) has product monodromy \(\gamma_j^\ell\)
on both sides. Thus the maps include the invariant stalks
and give the actual orthogonal parabolic decomposition
\[
 H^1(Y,j_*\mathcal M_Y)
     =H^1(\mathbb P^1,j_*\mathcal M)\perp
        H^1(\mathbb P^1,j_*(W\otimes\mathcal M)),
 \quad\dim_K H^1(Y,j_*\mathcal M_Y)=d_0+6.                \tag{S43.11}
\]
The first summand has its form multiplied by \(s\).

Here is the same comparison as an actual Jacobian projector.
For a nonzero \(u\in M(K)\), let \(e_u\in\mathbb Q[\Gamma]\)
be its orthogonal primitive coefficient idempotent, including
the conjugate representation. It is primitive because the
exact character field of \(\Gamma\) is \(K\).
Its actions commute with all the displayed quotient maps.
On \(JX\) put
\[
              \pi=e_u(1-s^{-1}f^*f_*),\qquad
                    A_X=\operatorname{im}(a\pi),          \tag{S43.12}
\]
where \(a>0\) is an integer clearing its coefficients.
Since \(f_*f^*=[s]\), this is an actual self-adjoint
rational projector and its indicated multiple is an
actual endomorphism. The containing primitive factor
\(R=\operatorname{im}(a e_u:JX\to JX)\) has dimension
\(6+d_0\). The fixed factor from \(JZ\) has dimension
\(d_0\), and \(A_X\) is its polarized orthogonal
complement in \(R\).

For \(w_s=se_1-\boldsymbol1\), its standard-group
primitive projector is exactly
\(e_s=\operatorname{av}_{L_0}-\operatorname{av}_{S_s}\).
Indeed the permutation representation is the direct sum of
its constant line and its standard module, and its
\(L_0\)-fixed standard line is \(Kw_s\); averaging
gives this literal identity. The tensor primitive corner is
\(e=e_s e_u\in\mathbb Q[H]\). Quotient pullback and
norm give the actual equality
\[
       |L_0|^{-1}\alpha^*\pi\alpha_*=e,
       \qquad
       (\alpha^*|_{A_X})^*(\theta_C|_{\operatorname{im}e})
                            =|L_0|\,\theta_X|_{A_X}.       \tag{S43.13}
\]
For the first identity use
\(\alpha^*\alpha_*=|L_0|\operatorname{av}_{L_0}\),
\((f\alpha)^*(f\alpha)_*=|S_s|\operatorname{av}_{S_s}\),
and \(|S_s|=s|L_0|\). It follows that
\(\alpha^*|_{A_X}\) is an actual isogeny onto the original
tensor primitive factor. All \(\Gamma\)-corner actions
commute with these maps. A different originally prescribed
primitive line is carried to \(Kw_s\otimes Ku\) by
the literal Fourier operator and polarization multiplier
(S41.15)--(S41.16). These equations identify the original
period map with that of the stated complementary factor.

\paragraph{5. The two period estimates.}
The first summand of (S43.11) is a fixed flat Hodge
subspace, after the finite label choice in paragraph3.
Pullback along the actual family of maps \(f\) is flat
on cohomology. Hence the original sixfold period map is
obtained from that of \(R\) by taking this fixed polarized
orthogonal complement. Its differential rank cannot exceed
that of \(R\). This is the actual direct-summand map of
Hodge filtrations, not a comparison of dimensions alone.

Let \(g=g(Y)\), and let \(b_Y\) be the number of
nontrivial coefficient branch values of the rank-two system
\(\mathcal M_Y\). Irreducibility and its dual give
\[
       6+d_0=2(2g-2)+\sum_{a=1}^{b_Y}c_a,
                   c_a\ge1,\qquad b_Y\le10+d_0-4g.         \tag{S43.14}
\]
The dimension of the pointed-base moduli is at most
\(3g-3+b_Y\le7+d_0-g\), using \(b_Y-3\) for genus
zero and \(b_Y\) for genus one. The unpointed genus-one
case is impossible here because (S43.14) has positive
left side. Finite monodromy cover choices over that pointed
base are discrete. Thus the actual family of primitive
factors \(R\) has period rank at most \(7+d_0\).
For \(d_0\le1\) this is at most eight, contrary to the
assumed rank nine of its complementary sixfold.

It remains to exclude \(d_0=2\). In this case (S43.8)
says that each of the three \(\gamma_j\) is fixed free.
It is also projectively nonidentity, so its two eigenvalues
are distinct roots of unity different from one. On a
permutation cycle of length \(\ell\), a fixed tensor
vector occurs once for each such eigenvalue whose
\(\ell\)-th power is one. There are no such eigenvalues
when \(\ell=1\); at most one when \(\ell=2\), since
only \(-1\) can occur; and at most two for \(\ell\ge3\).
In every case the fixed dimension is at most
\(2\ell/3\). Summing over the \(s\) letters gives
fixed dimension at most \(2s/3\) on \(K^s\otimes M\).
The constant \(M\)-summand is fixed free, so the same
fixed-space bound holds on the actual augmentation tensor.
The three inertia codimensions consequently satisfy
\[
    \sum_{j=1}^3c_j\ge3(2r-2s/3)=4r-2.
                                                               \tag{S43.15}
\]
At least nine further nonidentity inertia have codimension
at least two. Therefore
\(\sum c_j\ge4r-2+18=4r+16>4r+6\), the required
contradiction. This eigenvalue argument holds over every
coefficient field and uses no list of two-dimensional groups.

The multiplicity-two case is therefore excluded at period
rank nine for every imaginary quadratic \(K\). For the
Gaussian field and the seven other class-number-one fields
listed in HD-S42, combining this result with HD-S41,
HD-S42 and HD-G30 treats every nontrivial normal subgroup
generated by actual bireflections, without an upper bound
on coefficient rank. The remaining source case in that
scope has no such nontrivial subgroup; no assertion about
its classification or dominance is made here.
